\documentclass{article} % 假设你使用的是article类
\usepackage{tikz}
\usepackage{xcolor} % 确保添加了xcolor宏包
\usepackage{tkz-kiviat}
\begin{document}

\makeatletter
\def\tkz@KiviatGrad[#1](#2){% 
\begingroup
\pgfkeys{/kiviatgrad/.cd,
graduation distance= 0 pt,
prefix ={},
suffix={},
unity=1,
% added the following four lines
label precision/.store in=\gradlabel@precision,
label precision=1,
zero point/.store in=\tkz@grad@zero,
zero point=0
}
\pgfqkeys{/kiviatgrad}{#1}% 
\let\tikz@label@distance@tmp\tikz@label@distance
\global\let\tikz@label@distance\tkz@kiv@grad
  % started loop at zero, to get label for center point
 \foreach \nv in {0,...,\tkz@kiv@lattice}{
 % skipped the \pgfmathparse, moved
 % \tkz@kiv@unity*\nv-\tkz@grad@zero
 % into \pgfmathsetmacro. The last term is added
 \pgfmathsetmacro{\result}{\tkz@kiv@unity*\nv+20} % change from \pgfmathtruncatemacro
 % NOTE: use \gradlabel@precision for the precision
 \protected@edef\tkz@kiv@gd{%
    \tkz@kiv@prefix%
    \pgfmathprintnumber[precision=\gradlabel@precision,fixed]{\result}%% used \pgfmathprintnumber instead of "$\result$"
    \tkz@kiv@suffix} 
    \path[/kiviatgrad/.cd,#1] (0:0)--(360/\tkz@kiv@radial*#2:\nv*\tkz@kiv@gap)
       % NOTE: (360/\tkz@kiv@radial*#2)-90  -> (360/\tkz@kiv@radial*#2)
       % removed 90 because you rotate all the diagrams
       node[label=(360/\tkz@kiv@radial*#2):\tiny\tkz@kiv@gd] {}; % added \tiny
      }
 \let\tikz@label@distance\tikz@label@distance@tmp  
\endgroup
}%


\def\tkz@KiviatLine[#1](#2,#3){% 
\begingroup
\pgfkeys{/kiviatline/.cd,
fill= {},
opacity=.5,
% add these two lines, similar to for kiviatline
zero point/.store in=\tkz@line@zero,
zero point=0
}
%
% the code below has a lot of polar coordinates, of the form
%   (angle: <value> * <some factor>)
% all have been modified to get them on the form
%   (angle:{(<value>+\tkz@line@zero) * <some factor>})
% note extra braces, which are required when a component contains ()
% 
\pgfqkeys{/kiviatline}{#1}%   opacity ??????
\ifx\tkzutil@empty\tkz@kivl@fill \else 
\begin{scope}[on background layer]
 \path[fill=\tkz@kivl@fill,opacity=\tkz@kivl@opacity] (360/\tkz@kiv@radial*0:{(#2+\tkz@line@zero)*\tkz@kiv@gap*\tkz@kiv@step})   
\foreach \v [count=\rang from 1] in {#3}{%  
 -- (360/\tkz@kiv@radial*\rang:{(\v+\tkz@line@zero)*\tkz@kiv@gap*\tkz@kiv@step}) } -- (360/\tkz@kiv@radial*0:{(#2+\tkz@line@zero)*\tkz@kiv@gap*\tkz@kiv@step}); 
 \end{scope}
 \fi       
% added overlay option because something weird happened with the bounding box
\draw[#1,opacity=1,overlay] (0:{(#2+\tkz@line@zero)*\tkz@kiv@gap}) plot coordinates {(360/\tkz@kiv@radial*0:{(#2+\tkz@line@zero)*\tkz@kiv@gap*\tkz@kiv@step})}  
\foreach \v [count=\rang from 1] in {#3}{%  
 -- (360/\tkz@kiv@radial*\rang:{(\v+\tkz@line@zero)*\tkz@kiv@gap*\tkz@kiv@step}) plot coordinates {(360/\tkz@kiv@radial*\rang:{(\v+\tkz@line@zero)*\tkz@kiv@gap*\tkz@kiv@step})}} -- (360/\tkz@kiv@radial*0:{(#2+\tkz@line@zero)*\tkz@kiv@gap*\tkz@kiv@step});   
\endgroup
}%  


\makeatother
\definecolor{c1}{HTML}{81BECE} % 定义颜色 c1
\definecolor{c2}{HTML}{378BA4} % 定义颜色 c2
\definecolor{c3}{HTML}{FF5733} % 定义颜色 c3
\definecolor{c4}{HTML}{C70039} % 定义颜色 c4
\definecolor{c5}{HTML}{900C3F} % 定义颜色 c5
\definecolor{c6}{HTML}{581845} % 定义颜色 c6

% 定义全局缩放样式
\tikzset{global scale/.style={
    scale=#1,
    every node/.append style={scale=#1}
  }
}
\begin{tikzpicture}[
  label distance=13em, % 标签距离
  global scale = 0.73, % 全局缩放
  plot1/.style={ % 数据集1样式
    thin,
    draw=c1,
    mark=square*,
    mark options={
     ball color=c1,
     mark size=4pt
    },
    opacity=.5
  },
  plot2/.style={ % 数据集2样式
    thin,
    draw=c2,
    mark=triangle*,
    mark options={
      mark size=4pt,
      ball color=c2
    },
    opacity=.5
  },
  plot3/.style={ % 数据集3样式
    thin,
    draw=c3,
    mark=square*,
    mark options={
      mark size=4pt,
      ball color=c3
    },
    opacity=.5
  },
  plot4/.style={ % 数据集4样式
    thin,
    draw=c4,
    mark=triangle*,
    mark options={
      mark size=4pt,
      ball color=c4
    },
    opacity=.5
  },
  plot5/.style={ % 数据集5样式
    thin,
    draw=c5,
    mark=square*,
    mark options={
      mark size=4pt,
      ball color=c5
    },
    opacity=.5
  },
  plot6/.style={ % 数据集6样式
    thin,
    draw=c6,
    mark=triangle*,
    mark options={
      mark size=4pt,
      ball color=c6
    },
    opacity=.5
  }
]

\useasboundingbox (-25em,-5em) rectangle (10em,7em); % 设置边界框

\begin{scope}[scale=0.36,local bounding box=a,shift={(-50em,0)}] % 设置局部缩放和位置

\newcommand\KivStep{0.1}
\pgfmathsetmacro\Unity{1/\KivStep}
\newcommand\zeroshift{10.00}

 \tkzKiviatDiagram[
   radial  style/.style ={-},
   rotate=90, 
   lattice style/.style ={black!30},
   step=\KivStep,
   gap=1,
   lattice=7,  % 设置刻度数
]% 
{medqa,medmcqa,pubmedqa,careqa,AVG} % 定义雷达图标签

% 绘制数据集1的雷达图线条
\tkzKiviatLine[
  plot1
](44.46190102,46.56944776,47.3,52.03700409,47.59208822)

% 绘制数据集2的雷达图线条
\tkzKiviatLine[
  plot2
](0.628436764,0.262969161,1.9,0.56929372,0.840174911)

% 绘制数据集3的雷达图线条
\tkzKiviatLine[
  plot3
](10.44776119,5.426727229,5.3,11.68831169,8.215700028)

% 绘制数据集4的雷达图线条
\tkzKiviatLine[
  plot4
](47.28986646,49.43820225,45.9,59.95374489,50.6454534)

% 绘制数据集5的雷达图线条
\tkzKiviatLine[
  plot5
](49.48939513,50.60961033,43.7,61.21686533,51.2539677)

% 绘制数据集6的雷达图线条
\tkzKiviatLine[
  plot6
](47.60408484,50.10757829,42.6,60.46966732,50.19533261)

% 绘制刻度标签
\tkzKiviatGrad[unity=\Unity, label precision=0, zero point=\zeroshift](0)
\end{scope}
\end{tikzpicture}


\end{document}